Sigui $X$ una variable aleatòria que segueix una distribució binomial $B(6;\,0{,}8)$.

\begin{apartats}

\apartat[1,25]{0,75}
Calcula $P(X=6)$ i $P(X\ge5)$.

\begin{solucio}
$P(X=6)=0{,}8^6\approx0{,}262$ i $P(X=5)=6\cdot0{,}8^5\cdot0{,}2\approx0{,}393$, i per tant
$P(X\ge5)\approx0{,}655$.
\end{solucio}

\begin{tria}{binomial-tasca}
\itemtria{original}{1}{1,25}
Calcula $P(X>1)$ i $P(2\le X\le4)$.

\begin{solucio}
$P(X=0)=0{,}2^6=0{,}000064$ i $P(X=1)=6\cdot0{,}8\cdot0{,}2^5=0{,}001536$.\\
$P(X>1)=1-P(X\le1)=1-0{,}0016=0{,}9984$.\\
$P(2\le X\le4)=1-P(X\le1)-P(X\ge5)\approx1-0{,}0016-0{,}6554\approx0{,}343$.
\end{solucio}

\itemtria{inventa-context}{1}{1,25}
Inventa una situació real en què una variable segueixi aquesta distribució $B(6;\,0{,}8)$, i explica què
vol dir, en aquesta situació, el resultat de $P(X\ge5)$.

\begin{solucio}
Resposta oberta. Cal que hi hagi 6 proves independents, i a cada una, la mateixa probabilitat d'èxit,
$0{,}8$. Per exemple: una jugadora encerta el 80\,\% dels tirs lliures i en fa 6, i $X$ és el nombre
d'encerts.\\
Aleshores, $P(X\ge5)\approx0{,}655$ vol dir que, en un 65,5\,\% de les sèries de 6 tirs, n'encerta com a
mínim 5.
\end{solucio}
\end{tria}

\begin{nomesllarg}
\apartat{0,75}
Calcula la mitjana i la variància de $X$.

\begin{solucio}
$\mu=6\cdot0{,}8=4{,}8$ i $\sigma^2=6\cdot0{,}8\cdot0{,}2=0{,}96$.
\end{solucio}
\end{nomesllarg}

\end{apartats}
